\section{Study Discussion}
The low-to-moderate NASA-TLX scores suggest that the tool does not impose heavy cognitive demands, while the positive UEQ+ scores indicate a positive user experience across different dimensions. The custom questions help explain some results of the NASA-TLX and UEQ+ scores.

The results suggest that the flow graph fulfills the design goal of making the spreadsheet formula logic more accessible and more navigable. Another observation is that despite encountering specific usability problems, as mentioned by participants in the qualitative study, that the frustration ratings on the NASA-TLX remained low. Several participants encountered confusion when attemtping to directly edit formula output values in the flow graph, which the tool did not support. Yet, this did not have a big impact on the overall satisfaction, as seen in the UEQ+ Usefulness score and the general positive answers in the custom questions.
% TODO: refer to the study hermans detecting and visualizing inter-worksheet smells in spreadsheets => also came to the result that flow graph helper 

Nevertheless, the study also reveals ares where the design of the tool could be refined. The relative lower score, 1.45, of Dependability compared with the other scores, indicates that there is a gap between the users expected behavior and what the tool allows them to do, gulf of execution. %TODO: cite Paper of gulf of execution
In the study it was observed, that some users for example expected to be able to modify every value within the flow graph, but the system required them to navigate back to the worksheet and modify it there. This gap could be easily fixed by allowing the user to be able to change these values even though it potentially could destroy the current formula. Similarly, the feedback in the custom questions about the expansion point visibility and the interaction intensive task of expanding all reference nodes suggest shortcuts or more efficient traversal of the flow graph in deeply nested formulas. The next design iteration could for example introduce a collapse all or an expand all button.


Finally, all participants, regardless of prior spreadsheet knowledge, were able to complete the comprehension and manipulation tasks, which provides initial evidence for the tool's accessibility. Also, the perspicuity importance rating of 6.40, which was the highest amongst all other UEQ+ scores, underscores that for the users it is very important that the system is easy to learn. For future iterations it should be a design priority to maintain or even improve the learnability of the system, especially for users who may not understand function arguments or lookup semantics.
